
\documentclass{article}
\usepackage{colortbl}
\usepackage{makecell}
\usepackage{multirow}
\usepackage{supertabular}

\begin{document}

\newcounter{utterance}

\centering \large Interaction Transcript for game `eqbench', experiment `eq\_bench', episode 56 with reflexion{-}Qwen3.5{-}27B.
\vspace{24pt}

{ \footnotesize  \setcounter{utterance}{1}
\setlength{\tabcolsep}{0pt}
\begin{supertabular}{c@{$\;$}|p{.15\linewidth}@{}p{.15\linewidth}p{.15\linewidth}p{.15\linewidth}p{.15\linewidth}p{.15\linewidth}}
    \# & \multicolumn{2}{c}{Player} && \multicolumn{2}{c}{Game Master} \\
    \hline

    \theutterance \stepcounter{utterance}  
    & & & \multicolumn{4}{p{0.6\linewidth}}{
        \cellcolor[rgb]{0.9,0.9,0.9}{
            \makecell[{{p{\linewidth}}}]{
                \texttt{\tiny{[P1$\langle$GM]}}
                \texttt{You have attempted this same task before and did not succeed. These are the plans you wrote after those attempts.} \\
\texttt{Trial \#1: New plan: My previous attempt failed because I assigned a high score to "Excited" (6), which contradicts the character's defensive and sarcastic tone throughout the dialogue. Marcus is reacting to a personal jab about his reliability ("forgetting to show up") and immediately pivots to a competition to regain control, indicating he feels "Challenged" rather than genuinely happy or eager. In the future, I will strictly align the scores with the character's defensive posture: I will assign a very low score (0{-}1) to "Excited" to reflect his lack of genuine enthusiasm, keep "Humiliated" low as he deflects the insult, and ensure "Challenged" remains the dominant emotion with the highest score, accurately reflecting his acceptance of the bet as a counter{-}move.} \\
\\ 
\texttt{Trial \#2: New plan: My previous attempt failed because I assigned a non{-}zero score to "Excited" (1), which misinterprets Marcus's sarcastic and defensive tone as genuine enthusiasm. Marcus's reaction to the "low blow" about forgetting the party is to immediately issue a counter{-}challenge ("make this interesting," "Who can throw the better surprise"), which is a classic deflection mechanism to regain control, not an expression of excitement. In the future, I will strictly differentiate between "Challenged" (the primary emotion driving his competitive response) and "Excited" (which implies positive anticipation). I will assign a score of 0 to "Excited" to reflect his lack of genuine joy, keep "Humiliated" low (1{-}2) as he successfully deflects the insult, and ensure "Challenged" receives the highest score (9{-}10) to accurately capture his competitive pivot.} \\
\\ 
\\ 
\\ 
\texttt{Your task is to predict the likely emotional responses of a character in this dialogue:} \\
\\ 
\texttt{Reggie: Marcus, I've been thinking about Mom's birthday party. I've got some ideas.} \\
\texttt{Marcus: Oh, I bet you do. Another corporate gala, Reggie? } \\
\texttt{Reggie: Funny. I was actually thinking of something more... intimate. } \\
\texttt{Marcus: Intimate? You? Now, that's a surprise.} \\
\texttt{Reggie: Yes, well, I can be full of surprises. } \\
\texttt{Marcus: I'll believe it when I see it. } \\
\texttt{Reggie: What's that supposed to mean? } \\
\texttt{Marcus: Nothing. Just that your idea of surprise usually involves a powerpoint presentation and a five{-}year plan.} \\
\texttt{Reggie: And your idea of a surprise is forgetting to show up. } \\
\texttt{Marcus: Low blow, sis. } \\
\texttt{Reggie: Truth hurts, doesn't it? } \\
\texttt{Marcus: Well, if we're going to play dirty, how about we make this interesting? } \\
\texttt{Reggie: Interesting? } \\
\texttt{Marcus: Yeah, a little competition. Who can throw the better surprise for Mom. } \\
\texttt{Reggie: You're on. But don't say I didn't warn you when you're crying into your camera lens. } \\
\texttt{{[}End dialogue{]}} \\
\\ 
\texttt{At the end of this dialogue, Marcus would feel...} \\
\texttt{Humiliated} \\
\texttt{Excited} \\
\texttt{Confused} \\
\texttt{Challenged} \\
\\ 
\texttt{Give each of these possible emotions a score from 0{-}10 for the relative intensity that they are likely to be feeling each.} \\
\\ 
\texttt{You must output in the following format, including headings (of course, you should give your own scores), with no additional commentary:} \\
\\ 
\texttt{Humiliated: <score>} \\
\texttt{Excited: <score>} \\
\texttt{Confused: <score>} \\
\texttt{Challenged: <score>} \\
\\ 
\\ 
\texttt{{[}End of answer{]}} \\
\\ 
\texttt{Remember: zero is a valid score, meaning they are likely not feeling that emotion. You must score at least one emotion > 0.} \\
\\ 
\texttt{Your answer:} \\
            }
        }
    }
     \\ \\

    \theutterance \stepcounter{utterance}  
    & \multicolumn{4}{p{0.6\linewidth}}{
        \cellcolor[rgb]{0.561,0.737,0.561}{
            \makecell[{{p{\linewidth}}}]{
                \texttt{\tiny{[P1$\rangle$GM]}}
                \texttt{Humiliated: 2} \\
\texttt{Excited: 0} \\
\texttt{Confused: 0} \\
\texttt{Challenged: 9} \\
            }
        }
    }
    & & \\ \\

    \theutterance \stepcounter{utterance}  
    & & & \multicolumn{2}{p{0.3\linewidth}}{
        \cellcolor[rgb]{0.9,0.9,0.9}{
            \makecell[{{p{\linewidth}}}]{
                \texttt{\tiny{[GM$|$GM]}}
                \texttt{Humiliated: 2} \\
\texttt{Excited: 0} \\
\texttt{Confused: 0} \\
\texttt{Challenged: 9} \\
            }
        }
    }
    & & \\ \\

    \theutterance \stepcounter{utterance}  
    & & & \multicolumn{2}{p{0.3\linewidth}}{
        \cellcolor[rgb]{0.9,0.9,0.9}{
            \makecell[{{p{\linewidth}}}]{
                \texttt{\tiny{[GM$|$GM]}}
                \texttt{Humiliated: 0} \\
\texttt{Excited: 7} \\
\texttt{Confused: 0} \\
\texttt{Challenged: 8} \\
            }
        }
    }
    & & \\ \\

    \theutterance \stepcounter{utterance}  
    & & & \multicolumn{2}{p{0.3\linewidth}}{
        \cellcolor[rgb]{0.9,0.9,0.9}{
            \makecell[{{p{\linewidth}}}]{
                \texttt{\tiny{[GM$|$GM]}}
                \texttt{game\_result = LOSE} \\
            }
        }
    }
    & & \\ \\

    \theutterance \stepcounter{utterance}  
    & & \multicolumn{4}{p{0.6\linewidth}}{
        \cellcolor[rgb]{0.95,0.95,0.95}{
            \makecell[{{p{\linewidth}}}]{
                \texttt{\tiny{[P1$\leftrightarrow$P]}}
                \texttt{Post{-}episode Reflexion (not sent to GM):} \\
\texttt{New plan: My previous attempt failed because I assigned a non{-}zero score (2) to "Humiliated," which misinterprets Marcus's successful deflection of the insult. When Reggie calls out his unreliability ("Low blow, sis"), Marcus immediately pivots to a counter{-}challenge ("make this interesting") rather than showing signs of shame or defeat. This indicates he is not humiliated but rather combative. In the future, I will assign a score of 0 to "Humiliated" to reflect his lack of shame, keep "Excited" and "Confused" at 0 as they do not fit the defensive/competitive tone, and assign the maximum score (10) to "Challenged" to accurately capture his primary reaction of accepting the bet to regain control.} \\
            }
        }
    }
    & \\ \\

\end{supertabular}
}

\end{document}
